\documentclass[10pt,a4paper]{article}

\usepackage[T1]{fontenc}
\usepackage{lmodern}
\usepackage{amsmath,amssymb,mathtools}
\usepackage{booktabs,array}
\usepackage{geometry}
\usepackage{xcolor}
\usepackage{graphicx}
\usepackage{enumitem}
\usepackage{microtype}
\usepackage{fancyhdr}
\usepackage{hyperref}
\usepackage{listings}
\usepackage{caption}
\usepackage{float}

\geometry{left=1.75cm,right=1.75cm,top=1.65cm,bottom=1.65cm}
\setlength{\headheight}{14pt}
\setlength{\parindent}{0pt}
\setlength{\parskip}{0.35em}
\setlist[itemize]{leftmargin=1.5em,itemsep=0.15em,topsep=0.2em}
\setlist[enumerate]{leftmargin=1.7em,itemsep=0.15em,topsep=0.2em}

\definecolor{accent}{RGB}{29,66,105}
\definecolor{softblue}{RGB}{240,245,250}
\definecolor{codegray}{RGB}{247,247,247}
\hypersetup{
  colorlinks=true,
  linkcolor=accent,
  urlcolor=accent,
  citecolor=accent,
  pdftitle={A Source-Controlled Nonlinear Kapitza Model}
}
\captionsetup{font=small,labelfont=bf}
\lstset{
  basicstyle=\ttfamily\small,
  backgroundcolor=\color{codegray},
  frame=single,
  rulecolor=\color{black!15},
  columns=fullflexible,
  keepspaces=true,
  showstringspaces=false,
  breaklines=true
}

\pagestyle{fancy}
\fancyhf{}
\lhead{Source-controlled nonlinear Kapitza model}
\rhead{Analysis and reproducible computation}
\cfoot{\thepage}

\newcommand{\sech}{\operatorname{sech}}
\newcommand{\artanh}{\operatorname{artanh}}
\newcommand{\jumpu}{[u]}
\newcommand{\betahat}{\widehat{\beta}}
\newcommand{\qhat}{\widehat{q}}

\begin{document}

\begin{center}
  {\LARGE\bfseries\color{accent} A Source-Controlled Nonlinear Kapitza Model}\par\vspace{0.25em}
  {\Large Exact reduction, imperfect pitchfork unfolding, and S-curves}\par\vspace{0.6em}
  {\large Detailed report accompanying the Python and MATLAB implementations}
\end{center}

\begin{abstract}
This report develops a one-dimensional nonlinear Kapitza transmission model
in which the normalized strength \(\kappa\) of a positive volume heat source is
the bifurcation parameter. The underlying two-layer Poisson problem is the
one-dimensional non-uniqueness example in the recent work of Bejaoui, Ben
Belgacem, Jelassi, and Ma. Eliminating the two piecewise temperature fields
gives the jump relation
\(s=[u]=\kappa-(1+\delta)a\), where
\(\delta=\lambda/\mu\) is the ratio of the bulk conductivities. A specifically
constructed exponential contact conductance makes the scalar equilibrium
equation exactly equivalent to
\(x=\tanh(\kappa x+h)\). Consequently, the symmetric problem undergoes a
supercritical pitchfork at the critical source strength \(\kappa_c=1\).
Positive and negative biases unfold the pitchfork, the saddle-node trajectory
is analytic, and its local displacement satisfies a two-thirds power law.
The report also reconstructs the physical temperature fields, distinguishes
Picard-iteration stability from physical time-dependent stability, and
documents matching single-file Python and MATLAB implementations.
\end{abstract}

\tableofcontents

\section{The source problem from the two-layer Poisson model}

Consider the composite interval
\begin{equation}
  \Omega=[-1,1],\qquad
  \Omega_I=[-1,0],\qquad
  \Omega_E=[0,1],
\end{equation}
with a single interface at \(z=0\). Let \(\lambda>0\) and \(\mu>0\) be the
constant thermal conductivities in \(\Omega_I\) and \(\Omega_E\), respectively.
The endpoint temperatures are homogeneous,
\begin{equation}
  u_I(-1)=0,\qquad u_E(1)=0,
\end{equation}
and the positive volume source is supported only in the right layer:
\begin{equation}
  f(z)=2\mu\kappa\,\chi_{\Omega_E}(z),
  \qquad \kappa>0.
  \label{eq:source}
\end{equation}
The steady bulk equations are therefore
\begin{equation}
  -\lambda u_I''=0 \quad\hbox{in }\Omega_I,
  \qquad
  -\mu u_E''=2\mu\kappa \quad\hbox{in }\Omega_E.
  \label{eq:pde}
\end{equation}

This resolves an important notation issue. In this model \(\kappa\) is not a
bulk conductivity: the conductivities are \(\lambda\) and \(\mu\). After the
right equation in \eqref{eq:pde} is divided by \(\mu\), the normalized source
amplitude is \(2\kappa\). Thus a blackboard notation such as
\(\alpha=2\kappa\) should be read as a source-amplitude relation, not as the
old numerical sigmoid-slope convention. The general paper also uses the
letter \(\alpha\) for the piecewise conductivity field, so the symbols must be
kept separate.

\section{Exact elimination of the bulk temperature fields}

Solving \eqref{eq:pde} and imposing the exterior boundary conditions gives
\begin{equation}
  u(z)=u_{(a,b)}(z)=
  \begin{cases}
    a(z+1), & z\in\Omega_I,\\[0.2em]
    b(z-1)+\kappa(1-z^2), & z\in\Omega_E.
  \end{cases}
  \label{eq:piecewise}
\end{equation}
The interface traces are
\begin{equation}
  u_I(0)=a,
  \qquad
  u_E(0)=\kappa-b.
\end{equation}
Flux conservation and the nonlinear Kapitza law take the form
\begin{equation}
  \lambda a=\mu b
  =\beta_{\mathrm{phys}}(a,\kappa-b)(\kappa-b-a).
  \label{eq:transmission}
\end{equation}
Introduce the conductivity ratio
\begin{equation}
  \delta:=\frac{\lambda}{\mu}>0.
\end{equation}
Flux continuity gives \(b=\delta a\). With the jump convention
\(s=[u]=u_E(0)-u_I(0)\), one obtains the exact bulk relation
\begin{equation}
  \boxed{s=\kappa-(1+\delta)a.}
  \label{eq:jump}
\end{equation}
Define the conductance and flux normalized by \(\mu\),
\begin{equation}
  \betahat:=\frac{\beta_{\mathrm{phys}}}{\mu},
  \qquad
  \qhat:=\frac{q}{\mu}.
\end{equation}
Then \eqref{eq:transmission} becomes
\begin{equation}
  \boxed{\qhat=\delta a=\betahat(s)s.}
  \label{eq:balance}
\end{equation}
Solving \eqref{eq:jump}--\eqref{eq:balance} for \(a\) produces the scalar map
\begin{equation}
  a=F_{\kappa,h}(a):=
  \frac{\kappa\betahat_{\kappa,h}(s)}
       {\delta+(1+\delta)\betahat_{\kappa,h}(s)},
  \qquad s=\kappa-(1+\delta)a.
  \label{eq:Fgeneral}
\end{equation}

Because \(\betahat>0\), a positive steady heat flux gives \(a>0\) and
\(s>0\). Hence the physical interval moves with the source parameter:
\begin{equation}
  0<a<\gamma(\kappa),
  \qquad
  \gamma(\kappa):=\frac{\kappa}{1+\delta},
  \qquad
  0<s<\kappa.
  \label{eq:physical}
\end{equation}
The source term is essential. If both layers were source-free while the two
exterior temperatures were zero and the conductances were positive, the
standard energy identity would force the trivial solution.

\section{Exponential Kapitza law and normalized fixed-point equation}

The new constitutive choice is
\begin{equation}
  \boxed{
  \betahat_{\kappa,h}(s)
  =\frac{\delta}{1+\delta}
  \exp\!\left[2(\kappa-2s)+2h\right].}
  \label{eq:beta}
\end{equation}
It is positive, decreasing, and convex as a function of the jump:
\begin{equation}
  \partial_s\betahat_{\kappa,h}=-4\betahat_{\kappa,h}<0,
  \qquad
  \partial_{ss}\betahat_{\kappa,h}=16\betahat_{\kappa,h}>0.
\end{equation}
The bias has a direct constitutive interpretation:
\begin{equation}
  \betahat_{\kappa,h}(s)=e^{2h}\betahat_{\kappa,0}(s).
\end{equation}

Define the normalized state
\begin{equation}
  x:=\frac{2a}{\gamma(\kappa)}-1
  =\frac{2(1+\delta)a}{\kappa}-1
  =1-\frac{2s}{\kappa},
  \qquad -1<x<1.
  \label{eq:x}
\end{equation}
Writing \(E=\exp[2(\kappa-2s)+2h]\), equation \eqref{eq:Fgeneral} becomes
\begin{equation}
  F_{\kappa,h}(a)
  =\frac{\gamma(\kappa)}{2}
   \left[1+\tanh(\kappa-2s+h)\right].
\end{equation}
Since \(\kappa-2s=\kappa x\), the complete transmission problem is exactly
equivalent to
\begin{equation}
  \boxed{x=\tanh(\kappa x+h).}
  \label{eq:normalized}
\end{equation}
The material ratio \(\delta\) remains in the conversion back to \(a\), but it
drops out of the normalized root equation.

\begin{figure}[H]
  \centering
  \includegraphics[width=0.92\textwidth]{output/python/model_summary.png}
  \caption{The source, nonlinear conductance, scalar map, and heat-flux
  intersections are four views of the same reduced transmission problem.}
  \label{fig:summary}
\end{figure}

\section{Symmetric source-controlled pitchfork}

Set \(h=0\) and define
\(G(x)=\tanh(\kappa x)-x\). The central root \(x=0\) is always present and
\begin{equation}
  G'(0)=\kappa-1.
\end{equation}
It follows that
\begin{equation}
  \boxed{
  \begin{array}{ll}
    0<\kappa\leq1: & \text{one steady state},\\
    \kappa>1: & \text{three steady states}.
  \end{array}}
  \label{eq:count}
\end{equation}
Thus the critical source strength is
\begin{equation}
  \boxed{\kappa_c=1.}
\end{equation}

The central branch and its physical variables are
\begin{equation}
  x_c=0,
  \qquad
  a_c(\kappa)=\frac{\kappa}{2(1+\delta)},
  \qquad
  s_c(\kappa)=\frac{\kappa}{2}.
  \label{eq:center}
\end{equation}
Consequently, the central branch is horizontal in \((\kappa,x)\), but it is a
linearly increasing curve in the physical \((\kappa,a)\) plane. For
\(x\neq0\), the two outer branches have the exact parameterization
\begin{equation}
  \boxed{
  \kappa(x)=\frac{\artanh x}{x},
  \qquad
  a(x)=\frac{\kappa(x)(1+x)}{2(1+\delta)},
  \qquad 0<|x|<1.}
  \label{eq:outer}
\end{equation}
Since
\begin{equation}
  \frac{\artanh x}{x}=1+\frac{x^2}{3}+O(x^4),
\end{equation}
the outer states satisfy
\(x\sim\pm\sqrt{3(\kappa-1)}\): this is a supercritical pitchfork.

\begin{figure}[H]
  \centering
  \includegraphics[width=\textwidth]{output/python/fixed_point_regimes.png}
  \caption{Fixed-point intersections below, at, and above the critical source
  strength. Each panel uses its own physical interval
  \(0<a<\kappa/(1+\delta)\).}
  \label{fig:fixedpoints}
\end{figure}

\begin{figure}[H]
  \centering
  \includegraphics[width=0.92\textwidth]{output/python/symmetric_bifurcation.png}
  \caption{The standard pitchfork in normalized coordinates and the same
  branches reconstructed in the physical \((\kappa,a)\) plane.}
  \label{fig:pitchfork}
\end{figure}

\section{Picard stability and cobweb interpretation}

The numerical fixed-point iteration is
\begin{equation}
  x_{n+1}=\tanh(\kappa x_n+h).
\end{equation}
At a fixed point, its multiplier is
\begin{equation}
  m_{\mathrm{FP}}
  =\kappa\sech^2(\kappa x+h)
  =\kappa(1-x^2).
  \label{eq:multiplier}
\end{equation}
The notation \(m_{\mathrm{FP}}\) avoids confusion with the left conductivity
\(\lambda\). A root is locally attracting for this scalar iteration when
\(|m_{\mathrm{FP}}|<1\), and repelling when \(|m_{\mathrm{FP}}|>1\).

For \(h=0\), the central multiplier is \(m_{\mathrm{FP}}=\kappa\). Therefore
the central branch is attracting for \(\kappa<1\) and repelling for
\(\kappa>1\). The two outer branches are attracting. At \(\kappa=1\), the
linearization is neutral, but
\(\tanh x=x-x^3/3+O(x^5)\), so nearby iterations approach the origin only
algebraically. This produces the visible critical slowing in the cobweb plot.

\begin{figure}[H]
  \centering
  \includegraphics[width=\textwidth]{output/python/cobweb_regimes.png}
  \caption{Cobweb paths reveal convergence to the center below threshold,
  critical slowing at \(\kappa=1\), and selection of the two attracting outer
  branches above threshold.}
  \label{fig:cobweb}
\end{figure}

This classification concerns the Picard algorithm, equivalently the
discrete-time map used to compute the fixed points. It is not automatically a
proof of dynamical stability for a time-dependent heat equation; that would
require specifying and linearizing a transient transmission model.

\section{Nonzero bias and imperfect pitchfork unfolding}

For a fixed nonzero bias, solving \eqref{eq:normalized} for the control
parameter gives the exact branch parameterization
\begin{equation}
  \boxed{
  \kappa(x;h)=\frac{\artanh x-h}{x},
  \qquad
  a(x;h)=\frac{\kappa(x;h)(1+x)}{2(1+\delta)}.}
  \label{eq:biasedbranch}
\end{equation}
The symmetry is
\begin{equation}
  (x,h)\mapsto(-x,-h),
\end{equation}
so negative-bias diagrams are reflected versions of positive-bias diagrams in
normalized state space. In physical coordinates the moving factor
\(\gamma(\kappa)\) must still be applied.

\begin{figure}[H]
  \centering
  \includegraphics[width=\textwidth]{output/python/small_positive_bias_bifurcations.png}
  \caption{Six fixed positive biases with \(\kappa\) retained as the source
  control parameter. Solid branches are attracting for Picard iteration and
  dashed branches are repelling.}
  \label{fig:positive}
\end{figure}

\begin{figure}[H]
  \centering
  \includegraphics[width=\textwidth]{output/python/small_negative_bias_bifurcations.png}
  \caption{The corresponding six nonpositive biases. The preferred branch and
  the finite saddle node are reflected relative to the positive-bias case.}
  \label{fig:negative}
\end{figure}

\section{Saddle-node trajectory and local two-thirds law}

At a fold, the residual and its derivative vanish. Combining
\begin{equation}
  x=\tanh(\kappa x+h),
  \qquad
  1=\kappa(1-x^2),
\end{equation}
gives
\begin{equation}
  \boxed{
  \kappa_{\mathrm{SN}}=\frac{1}{1-x_{\mathrm{SN}}^2},
  \qquad
  h=\artanh(x_{\mathrm{SN}})
  -\frac{x_{\mathrm{SN}}}{1-x_{\mathrm{SN}}^2}.}
  \label{eq:fold}
\end{equation}
The physical critical state is
\begin{equation}
  a_{\mathrm{SN}}
  =\frac{\kappa_{\mathrm{SN}}(1+x_{\mathrm{SN}})}
         {2(1+\delta)}.
\end{equation}

Near the perfect pitchfork,
\begin{align}
  h&=-\frac{2}{3}x_{\mathrm{SN}}^3+O(x_{\mathrm{SN}}^5),\\
  \kappa_{\mathrm{SN}}-1
   &=x_{\mathrm{SN}}^2+O(x_{\mathrm{SN}}^4).
\end{align}
Eliminating \(x_{\mathrm{SN}}\) gives the local cusp scaling
\begin{equation}
  \boxed{
  \kappa_{\mathrm{SN}}-1
  \sim\left(\frac{3|h|}{2}\right)^{2/3}.}
  \label{eq:scaling}
\end{equation}

\begin{figure}[H]
  \centering
  \includegraphics[width=\textwidth]{output/python/critical_point_trajectory.png}
  \caption{The critical source strength, the physical fold state, and the
  agreement between the exact saddle-node trajectory and the local
  two-thirds law.}
  \label{fig:trajectory}
\end{figure}

\section{Fixed-source S-curve and heat-flux intersections}

For a fixed \(\kappa>1\), use \(h\) as the continuation parameter. The
S-curve has the exact physical parameterization
\begin{equation}
  \boxed{
  h(x)=\artanh x-\kappa x,
  \qquad
  a(x;\kappa)=\frac{\kappa(1+x)}{2(1+\delta)}.}
  \label{eq:scurve}
\end{equation}
The two fold states are
\begin{equation}
  x_f=\pm\sqrt{1-\frac{1}{\kappa}},
\end{equation}
with bias magnitude
\begin{equation}
  h_c=\kappa\sqrt{1-\frac{1}{\kappa}}
      -\artanh\!\sqrt{1-\frac{1}{\kappa}}.
\end{equation}

The same multiplicity can be displayed as heat-flux intersections. Eliminating
\(a\) from \eqref{eq:jump} gives the bulk-required normalized flux
\begin{equation}
  \qhat_{\mathrm{bulk}}(s;\kappa)
  =\frac{\delta}{1+\delta}(\kappa-s),
\end{equation}
whereas the interface law supplies
\begin{equation}
  \qhat_K(s;\kappa,h)=\betahat_{\kappa,h}(s)s.
\end{equation}
Their intersections are exactly the steady states.

\begin{figure}[H]
  \centering
  \includegraphics[width=\textwidth]{output/python/s_curve.png}
  \caption{For \(\delta=0.5\) and \(\kappa=1.5\), the fixed-source S-curve
  has two analytic saddle nodes. An interior fixed-bias slice has three
  equivalent heat-flux intersections.}
  \label{fig:scurve}
\end{figure}

\begin{figure}[H]
  \centering
  \includegraphics[width=\textwidth]{output/python/biased_fixed_point_cobweb.png}
  \caption{At a fixed interior bias, the three fixed-point intersections and
  two initial states on opposite sides of the repelling middle root display
  the two attracting basins.}
  \label{fig:biasedcobweb}
\end{figure}

\section{Reconstruction of the piecewise temperature fields}

Each scalar root gives a complete steady transmission solution through
\eqref{eq:piecewise}. For the representative symmetric case
\(\delta=0.5\), \(\kappa=1.5\), the normalized roots are approximately
\begin{equation}
  x=-0.8585596366,\qquad x=0,\qquad x=0.8585596366.
\end{equation}
The corresponding left interface traces are
\begin{equation}
  a\approx0.0707201817,\qquad0.5,\qquad0.9292798183.
\end{equation}
The discontinuities at \(z=0\) are not numerical artifacts: they are the
temperature jumps permitted by the nonlinear Kapitza contact resistance.

\begin{figure}[H]
  \centering
  \includegraphics[width=0.86\textwidth]{output/python/temperature_profiles.png}
  \caption{The three complete piecewise solutions reconstructed from the
  three scalar fixed points. Solid and dashed styles retain the Picard
  stability convention.}
  \label{fig:profiles}
\end{figure}

\section{Numerical implementation and reproducibility}

The Python and MATLAB programs implement the same workflow:
\begin{enumerate}
  \item solve \(x=\tanh(\kappa x+h)\) by exact symmetry when \(h=0\) and by
        robust bracketed root isolation when \(h\neq0\);
  \item reconstruct \(a,b,s,\betahat,\qhat\) and the multiplier from each root;
  \item parameterize the perfect and imperfect branches analytically;
  \item encode attracting segments as solid curves and repelling segments as
        dashed curves;
  \item generate matching PNG and vector PDF figures, plus CSV state tables.
\end{enumerate}

The Python entry point is
\begin{lstlisting}
python3 code/python/kapitza_kappa_model.py
\end{lstlisting}
when executed from the project directory. Optional arguments are
\texttt{--delta}, \texttt{--kappa}, and \texttt{--output}.

The MATLAB entry point is
\begin{lstlisting}
run('code/matlab/kapitza_kappa_model.m')
\end{lstlisting}
and requires no toolbox beyond standard MATLAB plotting and root finding. The
two programs are deliberately single-file implementations with English
comments throughout. The MATLAB script also accepts the environment variable
\texttt{KAPITZA\_STAGE=early} or \texttt{late} for systems that impose a short
batch-execution limit; without it, the complete workflow runs normally.

\section{Interpretation and scope}

The completed analysis separates three logically distinct components:
\begin{itemize}
  \item The two-layer Poisson equation, its volume source, the piecewise
        solution, and the scalar interface reduction come from the cited
        one-dimensional transmission example.
  \item The exponential conductance \eqref{eq:beta} is a new constitutive
        construction placed on that confirmed reduction. The source paper's
        principal bifurcation examples instead use power-law conductances.
  \item Stability in the present figures means local attraction or repulsion
        of the Picard fixed-point iteration. A claim about physical transient
        stability would require a separate time-dependent model.
\end{itemize}

Within this scope, \(\kappa\) is now a genuine physical control parameter: it
is the normalized amplitude of the positive source in \(\Omega_E\). The value
\(\kappa_c=1\) is a critical normalized source strength, not a condition that
the two material conductivities are equal. Material contrast is measured by
\(\delta=\lambda/\mu\).

\section*{Source}

E. Bejaoui, F. Ben Belgacem, F. Jelassi, and J. Ma,
\emph{Semilinear Elliptic Problems with NonLinear Contact Resistance in
Composite Domains}, HAL preprint, Section~7,
\href{https://hal.science/hal-05659943v1/document}
{hal-05659943v1}.

\end{document}
